% ---------------------------------------------------------------------
%  Chapter 16 : Coupled Fields, Dyson Dressing, the Operator IR, and the Bridge
% ---------------------------------------------------------------------
\chapter{Coupled Fields, Dyson Dressing, and the Bridge}\label{ch:spatial-coupled}

The two preceding spatial chapters built the machinery for the \emph{easy}
spatial case: a single physical field whose linear dynamics is a plain heat
kernel, $\Gret(t,k)=\Theta(t)\,\ee^{-(\mu+Dk^{2})t}$, and whose only
complication is a derivative-free interaction vertex. That covers a surprising
amount of physics, but it leaves three genuinely harder situations untouched,
and each one breaks a different assumption of the diagonal path:

\begin{enumerate}
  \item \term{Coupled multi-field theories.} When several fields
        $\phi_{1},\dots,\phi_{N}$ have linear dynamics that \emph{mix} --- an
        off-diagonal reaction (``mass'') matrix $M$ --- the free propagator is a
        \emph{matrix} exponential $\ee^{-Mt}$, not a product of independent
        scalar modes. The single-pole-per-field picture is simply wrong.
  \item \term{Unequal diffusion.} When the diffusion matrix $\mathcal{D}$ is not
        a scalar multiple of the identity, $M$ and $\mathcal{D}$ do not commute,
        so $\ee^{-(M+\mathcal{D}k^{2})t}$ does not factor and has no simple
        $k$-dependence. There is no closed-form propagator to read off.
  \item \term{Derivative vertices.} A vertex like Model~B's conserved
        $\Lap(\phi^{2})$, Burgers' $\partial_{x}(\phi^{2})$, or KPZ's
        $(\partial_{x}\phi)^{2}$ carries spatial derivatives. These cannot be
        folded into a momentum-independent coupling; each one deposits a
        \emph{momentum form factor} that the loop integrator must average.
\end{enumerate}

This chapter is the part of the engine that handles all three. It is, in
practice, four cooperating modules:
\file{spectral\_propagator.py} (the coupled free propagator),
\file{dyson\_dressing.py} (the unequal-diffusion correction series),
\file{spatial\_operator\_ir.py} (the symbolic operator algebra for derivative
vertices), and \file{pipeline\_bridge.py} (the 2157-line dispatcher that ties
everything to the shared diagram pipeline). A fifth module,
\file{loop\_dyson.py}, is an off-the-production-path cross-check oracle we
describe at the very end.

The organising idea --- stated in the bridge's own module docstring --- is
worth fixing in your mind before any of the math: \textbf{route through the
shared pipeline, then certify.} Rather than re-implementing diagrammatics for
the spatial case, the bridge substitutes $\Lap\to-q^{2}$ and runs the
\emph{same} enumeration, classification, and time-domain integration that a
time-only theory uses; it then \emph{certifies} that the diagram answer
$C(q,\tau)$ equals the per-mode structure read off the propagator, before doing
the external $q\to x$ Fourier transform analytically.

\begin{note}
Everything in this chapter sits \emph{behind} the public
\code{compute\_cumulants} call you met in the quickstart. You do not invoke
\code{compute\_coupled\_tree\_correlator} or \code{dressed\_tree\_C} by hand;
\code{compute\_cumulants} dispatches to them based on the shape of your theory
(coupled or not, equal or unequal diffusion, tree or loop, $k=2$ or $k\ge3$).
This chapter explains \emph{what} it dispatches to and \emph{why}.
\end{note}

\section{Where this subsystem sits}

The end-to-end flow, from the text you type to the array that comes out, runs:

\begin{lstlisting}[style=console]
theory file (.theory.py)
  -> SpatialTheoryBuilder / theory_compiler   (operator IR lowering lives here)
  -> FieldTheory.expand()                      (multivariate Taylor -> vertices)
  -> build_propagator / build_spatial_propagator   (the `prop` dict)
  -> compute_cumulants   (the public entry; dispatches by spatial_dim/coupling)
       -> compute_spatial_correlator_via_pipeline    (tree, diagonal)
       -> compute_coupled_tree_correlator            (tree, coupled)
       -> compute_spatial_correlator_generic         (loops, single-field)
       -> compute_coupled_loop_correlator            (loops, coupled)
       -> compute_spatial_kpoint / compute_coupled_kpoint   (k>=3)
  -> C(x, tau) array  ->  notebook / sim comparison
\end{lstlisting}

The five drivers in the middle are all in \file{pipeline\_bridge.py}, and the
chapter walks through them roughly in that order. The bridge \emph{feeds on} the
expanded \code{FieldTheory} object (\code{ft}), the model dict (\code{model}),
the propagator dict (\code{prop}), and the numeric parameters (\code{num\_params},
which crucially includes the mean-field saddle). It \emph{calls into} the shared
diagram pipeline, the spatial \code{full\_integrator} (the one genuine loop
integral, Chapter~\ref{ch:spatial-core}), the coupled-field modules of this
chapter, and the analytic $q\to x$ transforms in \code{spatial\_correlator}. It
\emph{produces} a real-space correlator array \code{C[len(tau\_grid),
len(spatial\_grid)]} plus an \code{info} dict of diagnostics.

\section{Why coupled fields need a different propagator}

Start with the linear theory, because the interaction structure rides on top of
it. For an $N$-component field the linearized inverse propagator --- the
bilinear ``kernel'' in the \msrjd{} action --- is, in the mixed
frequency/momentum representation $(\omega,k)$,
\begin{equation}
  K(\omega,k) \;=\; -\ii\omega\,I \;+\; M \;+\; \mathcal{D}\,|k|^{2}.
  \label{eq:Kwk}
\end{equation}
Two matrices appear, and both can be non-trivial:
\begin{itemize}
  \item $M$ is the \term{reaction matrix} (or mass matrix),
        $M=\mathrm{diag}(\mu_{i})-A^{(0)}$. It is the $k^{0}$
        (momentum-independent) part. Its off-diagonal entries are linear
        couplings between fields; diagonal $M$ means decoupled fields.
  \item $\mathcal{D}$ is the \term{diffusion matrix},
        $\mathcal{D}=\mathrm{diag}(D_{i})+A^{(2)}$. It is the $k^{2}$ part.
        $\mathcal{D}\propto I$ means equal (scalar) diffusion.
\end{itemize}
This is exactly the structure recorded in the \file{spectral\_propagator.py}
module docstring (lines~10--16), which names \eqref{eq:Kwk} and the two
splittings verbatim.

The retarded free propagator is the matrix Green's function
\begin{equation}
  \Gret(t,k)\;=\;\Theta(t)\,\ee^{-(M+\mathcal{D}|k|^{2})t}.
  \label{eq:GRmatrix}
\end{equation}
Here is the trouble in one sentence: \emph{if $M$ and $\mathcal{D}$ do not
commute, the exponent $M+\mathcal{D}|k|^{2}$ cannot be split, so
\eqref{eq:GRmatrix} is not a product of two matrix exponentials and has no
simple $k$-dependence.} The diagonal pipeline's whole strategy --- one
$\omega$-pole per field, $k$ entering only through $\mu+Dk^{2}$ --- collapses.

\subsection{The reference-diffusion split}

The escape is to split off a \emph{scalar} reference diffusion and treat the
rest as a perturbation. Write
\begin{equation}
  \mathcal{D}\;=\;D_{0}\,I\;+\;\widehat{\mathcal{D}},
  \qquad D_{0}\in\mathbb{R},\qquad
  \widehat{\mathcal{D}}=\text{residual}\ (=0\ \text{iff}\ \mathcal{D}\propto I).
  \label{eq:Dsplit}
\end{equation}
The reference kernel $K_{0}=-\ii\omega\,I+M+D_{0}|k|^{2}\,I$ now has scalar
diffusion, which commutes with \emph{everything}. The only remaining
non-commuting object is $M$ itself, and a diagonalizable $M$ can always be
resolved into its spectral projectors:
\begin{equation}
  M\;=\;\sum_{\alpha} m_{\alpha}\,P_{\alpha},
  \qquad \sum_{\alpha}P_{\alpha}=I,
  \qquad P_{\alpha}P_{\beta}=\delta_{\alpha\beta}P_{\alpha}.
  \label{eq:spectral}
\end{equation}
The $m_{\alpha}$ are the eigenvalues of $M$ and $P_{\alpha}$ is the rank-one
projector onto the $\alpha$-th eigenspace, built as the outer product of the
right eigenvector with the corresponding row of the inverse eigenvector matrix.
With this, the \term{reference propagator} is closed form (the paper's
Appendix~B, eq.~B23):
\begin{equation}
  G_{0}(t,k)\;=\;\Theta(t)\sum_{\alpha}P_{\alpha}\,
     \ee^{-(m_{\alpha}+D_{0}|k|^{2})t}
   \;=\;\Theta(t)\,\ee^{-Mt}\,\ee^{-D_{0}|k|^{2}t}.
  \label{eq:G0}
\end{equation}

Two exactness facts make this more than a formal manoeuvre, and both are stated
in the module docstring (lines~33--37) and validated in
\file{tests/test\_spectral\_propagator.py}:
\begin{itemize}
  \item \textbf{If $\widehat{\mathcal{D}}=0$} (scalar diffusion, even with a
        coupled $M$), then $G_{0}$ is the \emph{exact} full propagator. No
        series is needed. This alone unlocks every coupled-reaction,
        equal-diffusion theory.
  \item \textbf{If $M$ and $\mathcal{D}$ are both diagonal}, $G_{0}$ reduces to
        the per-field scalar heat kernel $\ee^{-(\mu_{i}+D_{i}|k|^{2})t}$ that
        the diagonal pipeline (\code{heat\_kernel.py}) already builds. So the
        coupled code degenerates gracefully to the case you already know.
\end{itemize}

\begin{defn}
A matrix $M$ is \term{diagonalizable} if it has a full set of linearly
independent eigenvectors, i.e. $M=V\,\mathrm{diag}(m_{\alpha})\,V^{-1}$ for an
invertible $V$. The spectral resolution \eqref{eq:spectral} exists exactly when
$M$ is diagonalizable. A \term{defective} matrix (repeated eigenvalues with
non-trivial Jordan blocks) has no such resolution and would need the resolvent /
confluent form; that case is deliberately deferred.
\end{defn}

\section{Building the projectors: \texttt{spectral\_propagator.py}}

The numeric core lives in five functions plus a cached dataclass. We take them
in the order the data flows.

\subsection{Splitting the diffusion}

\code{split\_reference\_diffusion(D\_mat, D0=None)}
(\file{spectral\_propagator.py:55}) implements \eqref{eq:Dsplit}:

\begin{lstlisting}
def split_reference_diffusion(D_mat, D0=None):
    D_mat = np.asarray(D_mat, dtype=float)
    n = D_mat.shape[0]
    if D0 is None:
        D0 = float(np.trace(D_mat) / n)
    Dhat = D_mat - D0 * np.eye(n)
    return float(D0), Dhat
\end{lstlisting}

When you do not pass a $D_{0}$, it defaults to the isotropic part
$\mathrm{trace}(\mathcal{D})/N$ --- the mean eigenvalue. That is a sensible
reference because it minimizes a natural norm of $\widehat{\mathcal{D}}$, which
(as \S\ref{sec:dyson} shows) is exactly what governs the convergence of the
Dyson series. The override exists so a caller can retune $D_{0}$ to shrink the
ratio $\lVert\widehat{\mathcal{D}}\rVert/D_{0}$.

\subsection{The spectral projectors and the defectiveness guard}

\code{spectral\_projectors(M)} (\file{spectral\_propagator.py:70}) is the
literal translation of \eqref{eq:spectral}:

\begin{lstlisting}
def spectral_projectors(M):
    M = np.asarray(M, dtype=complex)
    w, V = np.linalg.eig(M)
    cond = np.linalg.cond(V)
    if not np.isfinite(cond) or cond > _COND_CAP:
        raise ValueError(...)   # near-defective: deferred
    Vinv = np.linalg.inv(V)
    proj = [np.outer(V[:, a], Vinv[a, :]) for a in range(len(w))]
    return w, proj
\end{lstlisting}

Two NumPy calls do the algebra. \code{np.linalg.eig(M)} returns the eigenvalues
\code{w} (the $m_{\alpha}$) and the matrix \code{V} whose columns are the right
eigenvectors. The projector $P_{\alpha}=$ \code{np.outer(V[:,a], Vinv[a,:])} is
the outer product of the $\alpha$-th eigenvector with the $\alpha$-th row of
$V^{-1}$ --- the standard construction of the rank-one spectral projector of a
diagonalizable matrix.

\begin{defn}
\code{np.linalg.cond(V)} is the \term{condition number} of \code{V}: the ratio
of its largest to smallest singular value. It measures how close \code{V} is to
singular. A nearly defective $M$ has nearly parallel eigenvectors, so \code{V}
is nearly singular and \code{cond} blows up. The guard
\code{\_COND\_CAP = 1e10} (\file{spectral\_propagator.py:52}) is the line in the
sand: above it, the projector construction is numerically unreliable and the
function \emph{raises} \code{ValueError} rather than return garbage.
\end{defn}

\begin{gotcha}
The guard catches a \emph{near}-defective $M$ before it produces silently wrong
numbers, but it is the only protection: there is no fallback to a confluent or
resolvent form. If your theory has a genuinely defective reaction matrix (a
repeated eigenvalue with a Jordan block), this subsystem will refuse it. In
practice this is rare --- it requires a fine-tuned degeneracy --- but it is a
hard boundary of the implementation, not a soft warning.
\end{gotcha}

\subsection{The cached reference}

\code{SpectralReference} (\file{spectral\_propagator.py:91}) is a
\code{@dataclass} that caches the spectral data so it is computed once per
$(M,\mathcal{D})$ and reused across the whole $q\times\tau$ grid. It carries
\code{M}, \code{D} (the full diffusion), \code{D0}, \code{Dhat}, \code{eigvals}
(the $m_{\alpha}$), and \code{projectors} (the $P_{\alpha}$). Three properties
encode the facts above:

\begin{itemize}
  \item \code{n\_fields} $=$ \code{M.shape[0]}.
  \item \code{is\_scalar\_diffusion} $=$ \code{bool(np.allclose(self.Dhat, 0.0))}.
        \textbf{When this is true, $G_{0}$ is the exact full propagator and no
        Dyson series is needed.} This single boolean is the most consequential
        branch point in the coupled tree path.
  \item \code{G0(ksq, t)} returns the matrix
        $\sum_{\alpha}P_{\alpha}\,\ee^{-(m_{\alpha}+D_{0}\cdot\code{ksq})t}$
        (the caller applies $\Theta(t)$).
\end{itemize}

\code{build\_reference(M, D, D0=None)} (\file{spectral\_propagator.py:116}) is
the single constructor every coupled path uses: it validates that $M$ and $D$
are square and the same shape, splits the diffusion, computes the projectors,
and assembles the dataclass.

\begin{note}
A \code{@dataclass} is a Python class whose boilerplate (the constructor, a
readable \code{repr}, equality) is generated automatically from a few typed
field declarations. You write the field list; Python writes
\code{\_\_init\_\_}. It is the idiomatic way to carry a small bundle of related
values, and you will see it again for the integrator's diagram descriptors.
\end{note}

\section{The tree-level coupled two-point}

The free two-point of a coupled \emph{linear} theory is \emph{not} a sum of
independent Ornstein--Uhlenbeck modes --- that diagonal form is precisely what
the single-field path builds and what fails here. The correct object comes from
two classical results of linear stochastic dynamics.

Let $A(q)=M+\mathcal{D}|q|^{2}$ be the relaxation matrix at momentum $q$, and
let $N$ be the symmetric noise covariance, read from the $(2,0)$ response-field
sector of the action (the statement $\avg{\xi\xi^{\mathsf T}}=N$). Then:

\begin{description}
  \item[Lyapunov / stationary covariance.] The equal-time covariance
        $\Sigma(q)$ solves the \term{continuous Lyapunov equation}
        \begin{equation}
          A(q)\,\Sigma(q)\;+\;\Sigma(q)\,A(q)^{\mathsf T}\;=\;N.
          \label{eq:lyap}
        \end{equation}
  \item[Fluctuation--regression.] The relaxation of an equilibrium fluctuation
        follows the same exponential law as a macroscopic perturbation
        (Onsager's regression hypothesis), so
        \begin{equation}
          C(q,\tau)=\ee^{-A(q)|\tau|}\,\Sigma(q)\ \ (\tau\ge0),
          \qquad
          C(q,-\tau)=\Sigma(q)\,\ee^{-A(q)^{\mathsf T}|\tau|},
          \label{eq:fdt}
        \end{equation}
        i.e. $C(q,\tau)=C(q,|\tau|)^{\mathsf T}$.
\end{description}

For scalar diffusion $A(q)=M+D_{0}|q|^{2}\,I$ shares $M$'s eigenprojectors, so
$\ee^{-A|\tau|}$ \emph{is} $G_{0}(|q|^{2},|\tau|)$ and \eqref{eq:fdt} is exact.
Unequal diffusion needs the Dyson series of the next section.

\subsection{\texttt{coupled\_two\_point} and the Lyapunov solver}

\code{lyapunov\_covariance(A, N)} (\file{spectral\_propagator.py:156}) is a
single call to SciPy:

\begin{lstlisting}
def lyapunov_covariance(A, N):
    A = np.asarray(A, dtype=float)
    N = np.asarray(N, dtype=float)
    return solve_continuous_lyapunov(A, N)
\end{lstlisting}

\code{coupled\_two\_point(ref, N, qsq, tau)}
(\file{spectral\_propagator.py:169}) assembles \eqref{eq:lyap}--\eqref{eq:fdt}:

\begin{lstlisting}
def coupled_two_point(ref, N, qsq, tau):
    if not ref.is_scalar_diffusion:
        raise ValueError(...)          # unequal diffusion needs Dyson
    n = ref.n_fields
    A = np.asarray(ref.M, dtype=float) + ref.D0 * float(qsq) * np.eye(n)
    Sigma = lyapunov_covariance(A, N)
    G = expm(-A * abs(float(tau)))     # e^{-A|tau|}  (= G0 for scalar D)
    return G @ Sigma if tau >= 0 else Sigma @ G.T
\end{lstlisting}

It forms $A=M+D_{0}q^{2}I$, solves for $\Sigma$, exponentiates, and returns
$G\Sigma$ (or its transpose for $\tau<0$). The returned $(N,N)$ matrix has
$C_{ij}$ equal to the cross-correlation of fields $i$ and $j$. The
\code{raise} on the first line is the guard: this routine is for scalar
diffusion only.

\begin{defn}
\code{scipy.linalg.solve\_continuous\_lyapunov(A, N)} solves
$A\,X+X\,A^{\mathsf H}=Q$ for $X$ (here $Q=N$ and $A^{\mathsf H}=A^{\mathsf T}$
since $A$ is real). It uses the \term{Bartels--Stewart algorithm} (a Schur
decomposition followed by back-substitution), which is far more stable than
forming and inverting the $N^{2}\times N^{2}$ Kronecker-product system by hand.
\code{scipy.linalg.expm(Z)} is the \term{matrix exponential} $\ee^{Z}$, computed
by a Pad\'e approximation with scaling-and-squaring --- robust even for
non-normal matrices, which is essential here because $A$ need not be symmetric.
\end{defn}

\section{The Dyson--Duhamel series for unequal diffusion}\label{sec:dyson}

When $\widehat{\mathcal{D}}\ne0$ there is no closed-form propagator, and we
expand $\Gret$ in powers of the insertion $\widehat{\mathcal{D}}|k|^{2}$ (paper
Appendix~B, \S B.24--B.30):
\begin{align}
  \Gret(t,k)&=\sum_{n\ge0}G_{n}(t,k),\notag\\
  G_{n}(t,k)&=(-|k|^{2})^{n}\,\ee^{-D_{0}|k|^{2}t}\,\mathcal{H}_{n}(t).
  \label{eq:B26}
\end{align}
The matrix factor $\mathcal{H}_{n}(t)$ comes from the $n$-fold \term{Duhamel
convolution} --- each $\widehat{\mathcal{D}}$ insertion sandwiched between
reference evolutions:
\begin{equation}
  G_{n}(t)=(-|k|^{2})^{n}\!\!\int\limits_{t\ge s_{1}\ge\cdots\ge s_{n}\ge0}
   \!\!\ee^{-M(t-s_{1})}\widehat{\mathcal{D}}\,
   \ee^{-M(s_{1}-s_{2})}\widehat{\mathcal{D}}\cdots
   \ee^{-Ms_{n}}\,\dd\mathbf{s}.
  \label{eq:duhamel}
\end{equation}
Inserting the spectral resolution \eqref{eq:spectral} for every $\ee^{-M(\cdot)}$
turns this nested time integral into a \term{divided difference}, and
shift-invariance factors out $\ee^{-m_{\alpha_{0}}t}$:
\begin{equation}
  \mathcal{H}_{n}(t)=\!\!\sum_{\alpha_{0}\dots\alpha_{n}}\!\!
   P_{\alpha_{0}}\widehat{\mathcal{D}}P_{\alpha_{1}}\cdots
   \widehat{\mathcal{D}}P_{\alpha_{n}}\,
   \ee^{-m_{\alpha_{0}}t}\,
   \Phi_{n}\!\big(t;\,m_{\alpha_{1}}-m_{\alpha_{0}},\dots,
   m_{\alpha_{n}}-m_{\alpha_{0}}\big).
  \label{eq:B27}
\end{equation}

\subsection{The one new primitive: \texorpdfstring{$\Phi_{n}$}{Phi\_n}}

Everything in \eqref{eq:B27} is projector bookkeeping except the scalar function
$\Phi_{n}$, the nested-simplex time integral
\begin{equation}
  \Phi_{n}(t;\nu_{1},\dots,\nu_{n})=\int_{\sigma_{n}}
    t^{n}\,\ee^{-t\sum_{i}u_{i}\nu_{i}}\,\dd\mathbf{u},
  \qquad \Phi_{0}(t)=1,
  \label{eq:phin}
\end{equation}
over the standard simplex $\sigma_{n}=\{u_{i}\ge0,\ \sum u_{i}\le1\}$. The
classical \term{Hermite--Genocchi formula} with $f(z)=\ee^{-tz}$ (so
$f^{(n)}(z)=(-t)^{n}\ee^{-tz}$) and nodes $\{0,\nu_{1},\dots,\nu_{n}\}$ states
that this simplex integral is, up to sign, the $n$-th divided difference of $f$:
\begin{equation}
  \Phi_{n}(t;\nu)=(-1)^{n}\,f[0,\nu_{1},\dots,\nu_{n}].
  \label{eq:hg}
\end{equation}

\begin{defn}
The \term{divided difference} $f[x_{0},\dots,x_{n}]$ is the leading coefficient
of the degree-$n$ polynomial interpolating $f$ at the nodes $x_{0},\dots,x_{n}$;
for distinct nodes it equals $\sum_{i}f(x_{i})/\prod_{j\ne i}(x_{i}-x_{j})$. The
\term{Hermite--Genocchi formula} writes that same object as a simplex integral
of $f^{(n)}$, which is what links the Duhamel integral \eqref{eq:duhamel} to the
divided difference.
\end{defn}

\subsection{Confluent-safe evaluation: the Opitz theorem}

The naive formula $\sum_{i}f(x_{i})/\prod_{j\ne i}(x_{i}-x_{j})$ is a numerical
trap: it is $0/0$ whenever two nodes coincide or nearly coincide. And the
eigenvalue differences $\nu=m_{\alpha_{i}}-m_{\alpha_{0}}$ \emph{routinely}
repeat --- whenever several $m_{\alpha}$ are equal, or come as the
complex-conjugate pairs that a real $M$ produces. The \term{Opitz theorem}
sidesteps division entirely:

\begin{quote}
For analytic $f$, build the $(n+1)\times(n+1)$ upper-bidiagonal matrix $Z$ with
the nodes $\{0,\nu_{1},\dots,\nu_{n}\}$ on the diagonal and ones on the
superdiagonal. Then $f(Z)[0,n]=f[0,\nu_{1},\dots,\nu_{n}]$ --- the top-right
corner of $f$ applied to $Z$.
\end{quote}

So $\Phi_{n}(t;\nu)=(-1)^{n}\,\code{expm}(-t\,Z)[0,n]$: a single matrix
exponential of a tiny $(n+1)\times(n+1)$ matrix, exact at coincident nodes and
stable at near-coincident ones. The implementation
(\file{spectral\_propagator.py:221}) is:

\begin{lstlisting}
def phi_n(t, nus):
    nus = np.asarray(nus, dtype=complex).ravel()
    n = nus.size
    if n == 0:
        return complex(1.0)
    Z = _opitz_bidiagonal(nus)
    return complex((-1) ** n * expm(-float(t) * Z)[0, n])
\end{lstlisting}

with \code{\_opitz\_bidiagonal} (\file{spectral\_propagator.py:211}) building
$Z$ via \code{np.diag(np.ones(n), k=1)} (ones on the superdiagonal) and then
writing the nodes onto the diagonal.

\begin{example}\label{ex:phi1}
The $n=1$ check from the docstring (\file{spectral\_propagator.py:240}). With
$Z=\begin{psmallmatrix}0&1\\0&\nu\end{psmallmatrix}$,
$\code{expm}(-tZ)[0,1]=-t\,(1-\ee^{-t\nu})/(t\nu)=-(1-\ee^{-t\nu})/\nu$; times
$(-1)$ gives $\Phi_{1}=(1-\ee^{-t\nu})/\nu$, which is exactly
$\int_{0}^{1}t\,\ee^{-tu\nu}\,\dd u$. The formula reduces to the textbook answer,
and --- crucially --- it has a finite limit $\Phi_{1}\to t$ as $\nu\to0$ with no
cancellation, because the matrix exponential never divides by $\nu$.
\end{example}

\begin{gotcha}
\code{phi\_n} always returns a Python \code{complex}, never a float, even when
the inputs look real. This is deliberate: a real $M$ can have a complex
spectrum, so the node differences $\nu$ come in conjugate pairs. The callers
decide when it is safe to take the real part. \code{phi\_n\_batch}
(\file{spectral\_propagator.py:259}) is the vectorized version --- one
\code{expm} per $t$ on the shared tiny $Z$ --- used by the dressing assembly.
\end{gotcha}

\subsection{Assembling the dressing: \texttt{dyson\_dressing.py}}

\code{hcal\_n(ts, M, Dhat, n)} (\file{dyson\_dressing.py:47}) is the direct
transcription of \eqref{eq:B27}:

\begin{lstlisting}
def hcal_n(ts, M, Dhat, n):
    eig, proj = spectral_projectors(M)
    nf = len(eig)
    out = np.zeros((ts.size, nf, nf), dtype=complex)
    for string in itertools.product(range(nf), repeat=n + 1):
        mat = proj[string[0]]
        for a_i in string[1:]:
            mat = mat @ Dhat @ proj[a_i]
        if not np.any(np.abs(mat) > 1e-300):
            continue                                   # skip a zero string
        m0 = eig[string[0]]
        nus = [eig[a_i] - m0 for a_i in string[1:]]
        scal = np.exp(-m0 * ts) * phi_n_batch(ts, nus)
        out += scal[:, None, None] * mat[None, :, :]
    return out
\end{lstlisting}

It loops over every projector string $(\alpha_{0},\dots,\alpha_{n})$ ---
\code{itertools.product(range(nf), repeat=n+1)} enumerates all $N^{n+1}$ of
them --- builds the outer-product chain
$P_{\alpha_{0}}\widehat{\mathcal{D}}P_{\alpha_{1}}\cdots\widehat{\mathcal{D}}P_{\alpha_{n}}$,
skips strings whose matrix is numerically zero, forms the scalar
$\ee^{-m_{\alpha_{0}}t}\Phi_{n}$ via \code{phi\_n\_batch}, and accumulates the
$\text{scalar}\otimes\text{matrix}$ product. The cost is $N^{n+1}$ strings ---
acceptable because the truncation order $n$ stays small ($\lesssim4$).

\begin{defn}
\code{itertools.product(range(nf), repeat=m)} is the Python standard-library
Cartesian-product generator: it yields every length-$m$ tuple of values drawn
from $0,\dots,nf-1$, i.e. exactly the index strings summed in \eqref{eq:B27}.
\code{np.einsum} (used below) is NumPy's Einstein-summation engine: you write the
index pattern of a tensor contraction as a string and it performs the sum in
compiled C.
\end{defn}

\code{dressed\_GR(ts, qsq, M, Dhat, D0, order)} (\file{dyson\_dressing.py:73})
assembles \eqref{eq:B26}: $\sum_{n=0}^{\text{order}}(-q^{2})^{n}\mathcal{H}_{n}$,
multiplied by $\ee^{-D_{0}q^{2}t}$. At \code{order=0} it is exactly the scalar
reference $G_{0}$.

\subsection{The dressed tree two-point and its convergence}

With the dressed retarded propagator in hand, the dressed tree two-point is the
same $\int\,\dd s\,\Gret N\Gret^{\mathsf T}$ structure as
\eqref{eq:fdt}, now evaluated by quadrature because $\Gret$ is a series rather
than a single exponential:
\begin{equation}
  C^{(N)}(q,\tau\ge0)=\int_{0}^{\infty}\dd s\;
    \Gret^{(N)}(\tau+s,q)\,N\,\Gret^{(N)}(s,q)^{\mathsf T}.
  \label{eq:dressedC}
\end{equation}
\code{dressed\_tree\_C(q, tau, M, Dhat, D0, N, order, n\_s=48,
s\_cap\_scale=32.0)} (\file{dyson\_dressing.py:87}) does this:

\begin{lstlisting}
    cap = s_cap_scale / (mu_min + float(D0) * float(q) ** 2)
    xv, wv = np.polynomial.legendre.leggauss(n_s)
    vv = 0.5 * (xv + 1.0)
    s_nodes = cap * vv * vv
    s_w = wv * cap * vv                                  # int ds = sum w*cap*v
    ...
    C = np.einsum('s,sij,jk,slk->il', s_w, G_late, N, G_early)
\end{lstlisting}

Two ideas are worth unpacking. First, it validates that
$\min\mathrm{Re}\,\mathrm{eig}(M)>0$ (a stable theory) and \emph{raises}
otherwise. Second, the $s$-integral uses a \term{$\sigma$-concentrated
substitution} $s=\code{cap}\cdot v^{2}$: \term{Gauss--Legendre quadrature} (fixed
nodes \code{xv} and weights \code{wv} on $[-1,1]$, from \code{leggauss}) is mapped
through $s=\code{cap}\,v^{2}$ so the nodes pack near $s=0$, where the integrand is
largest, with \code{cap} set by the slowest eigenvalue. This is the same trick
the chamber integrator uses, and it makes a modest \code{n\_s=48} accurate.

\begin{defn}
\term{Gauss--Legendre quadrature} approximates $\int_{-1}^{1}g\,\dd v$ by a
weighted sum $\sum_{i}w_{i}\,g(v_{i})$ at fixed nodes; an $n$-point rule is exact
for polynomials of degree $\le 2n-1$. \code{np.polynomial.legendre.leggauss(n)}
returns the nodes and weights. The substitution $s=\code{cap}\,v^{2}$ turns the
half-line $\int_{0}^{\infty}$ into $\int_{0}^{1}$ while concentrating sample
points near the origin.
\end{defn}

The convergence of the truncation is governed by the \term{residual ratio}
\begin{equation}
  \rho=\frac{\lVert\widehat{\mathcal{D}}\rVert}{D_{0}}
       =\frac{\max|\mathrm{eig}(\widehat{\mathcal{D}})|}{D_{0}}.
  \label{eq:rho}
\end{equation}
The series is geometric in the order at large $q$ and \emph{exact at $q=0$}
(where the insertion $\widehat{\mathcal{D}}|k|^{2}$ vanishes), so convergence
requires $\rho<1$. The bridge checks this (\S\ref{sec:dispatch}) and you can
shrink $\rho$ by retuning $D_{0}$.

\section{Loop-level coupled fields: spectral assignments}

So far this has all been tree-level (the linear theory). At loop order the
interaction vertices return, and a beautiful simplification appears for scalar
diffusion.

With $\widehat{\mathcal{D}}=0$, \emph{every} edge's momentum factor is the same
heat kernel $\ee^{-D_{0}k^{2}w}$. The Symanzik (parametric-momentum) machinery
of Chapter~\ref{ch:spatial-core} --- which only ever sees scalar diffusion --- is
therefore completely untouched. The coupling lives entirely in the time/matrix
factor: each retarded segment carries $\sum_{\alpha}P_{\alpha}\,\ee^{-m_{\alpha}w}$.
Expanding every segment (a retarded ``R'' edge is one segment; a correlation
``C'' edge is two glued half-segments) turns one coupled diagram into a
\emph{weighted sum of scalar diagrams}:
\begin{equation}
  \text{value}(\Gamma)=\code{pv}\cdot
   \sum_{\{\alpha_{r}\}}\prod_{r}[P_{\alpha_{r}}]_{p_{r}r_{r}}\cdot
   I_{\text{spec}}\big(\{m_{\alpha_{r}}\}\big).
  \label{eq:specassign}
\end{equation}
Here $(r_{r},p_{r})$ are the segment's (response, physical) matrix indices,
threaded through the edge descriptor's \code{fpairs}; \code{pv} is the
enumeration's $\Sym(\Gamma)\cdot\text{prefactor}$; and $I_{\text{spec}}$ is the
shared Symanzik/chamber quadrature evaluated for one mass assignment, with all
assignments batched against a single integral pass.

\code{compute\_coupled\_loop\_correlator} (\file{pipeline\_bridge.py:1090})
implements \eqref{eq:specassign}. Its docstring (lines~1104--1110) states the
formula and one subtle normalization point that is worth flagging:

\begin{gotcha}
The single-field generic path multiplies by a universal $2^{-n_{C}}$ factor (the
one-segment ``C'' convention). The coupled spectral-assignment paths do
\emph{not}: the two-glued-segment ``C'' representation natively produces the
Lyapunov denominator $1/(m_{\alpha}+m_{\beta}+2D_{0}k^{2})$, so the conversion
does not apply (\file{pipeline\_bridge.py:1108-1110}). Mixing the two conventions
is a classic source of a factor-of-$2^{n_{C}}$ error.
\end{gotcha}

Mechanically, the driver maps each enumerated loop diagram to the integrator's
``C-stack'' descriptor (\code{diagram\_to\_cstack}), expands it into rows
(\code{spectral\_rows}), reads the $(\text{resp},\text{phys})$ indices off each
edge's \code{fpairs}, and builds the full assignment grid with
\code{itertools.product(range(nf), repeat=n\_rows)}. Per-assignment weights are
the product $\prod_{r}[P_{\alpha_{r}}]_{p_{r}r_{r}}$, the live ones are kept, and
the mass table $\{m_{\alpha_{r}}\}$ feeds the one shared quadrature pass.

\begin{gotcha}
\textbf{Cross-correlator $\tau$-asymmetry (the $i\ne j$ fix).} For an
\emph{auto}-correlator ($i=j$) the two mirror-image records dedupe to one record
plus the $\Gamma(\tau)+\Gamma(-\tau)$ retarded completion. For a
\emph{cross}-correlator ($i\ne j$) \emph{both} mirror records survive and each is
evaluated once at the times matching its own leaf field order --- applying the
$\pm\tau$ completion there would double-count and symmetrize away the genuine
$\tau$-asymmetry of $C_{ij}$ (\file{pipeline\_bridge.py:1368-1400}). The
coupled-loop and coupled-$k$-point drivers handle this with the orbit-stabilizer
mapping sum (\S\ref{sec:kpoint}); the comment notes the two views coincide at
$k=2$.
\end{gotcha}

When the diffusion is unequal, the loop path layers the Dyson $\mathcal{H}_{n}$
insertions onto individual segments via a generalized partial-fraction expansion
of the $n$-fold Duhamel convolution (the \code{\_pf\_labels} / \code{\_hn\_labels}
machinery, with a log-derivative Bell expansion for repeated masses), and refuses
periodic boundaries or derivative vertices in this v1.

\section{The operator IR for derivative vertices}

We now switch from the matrix structure of coupled fields to the third hard
case: vertices that carry spatial derivatives. The home for these is
\file{spatial\_operator\_ir.py}, a small \term{intermediate representation} (IR)
that hosts the operators $\Lap$ ($\nabla^{2}$), $\partial_{t}$, and
$\partial_{x_{i}}$ as symbolic nodes and gives them an explicit algebra.

\subsection{Why an IR, and why inert nodes}

The module docstring (lines~10--20) argues the design from two failed
alternatives. A bare multiplicative symbol --- writing
\code{D*Laplacian*phi} --- loses \emph{which field the derivative acts on}
(\code{phi*Laplacian*phi} is ambiguous) and carries no operator algebra, so
linearity and the homogeneous-saddle annihilation would have to be hand-patched
downstream. Sage's own \code{laplacian()} (the Laplace--Beltrami operator from
SageManifolds) wants a manifold, a metric, and \code{DiffScalarField} objects,
and does not compose with the \msrjd{} field-doubling and multivariate-Taylor
machinery, which works on commutative polynomial-ring generators.

The solution is to host the operators as \term{inert Sage function
applications}:

\begin{lstlisting}
from sage.all import SR, I, function
_LAP = function('Lap')              # nabla^2 (.)
_DT  = function('Dt')               # d_t (.)
_DX  = function('Dx')               # d_{x_i} (., i)
\end{lstlisting}

(\file{spatial\_operator\_ir.py:59,68-70}).

\begin{defn}
Sage's \code{function('Lap')} creates an \term{inert symbolic function}: a
callable symbol with no defined evaluation rule. So \code{Lap(phi)} is just a
tree node \code{Lap(phi)} that survives substitution, printing, and tree-walking
unchanged --- Sage never tries to ``compute'' it. The entire IR design rests on
this: the operators are inert nodes, and \emph{all} semantics live in explicit
passes. Note also that Sage's \term{Symbolic Ring} \code{SR} caches symbols by
name, so \code{SR.var('Laplacian')} anywhere returns the same object --- a
property the bridge exploits when it substitutes $\Lap\to-q^{2}$.
\end{defn}

A companion bare symbol \code{GRADX\_SYM = SR.var('GradX', ...)}
(\file{spatial\_operator\_ir.py:78}) is what a \emph{bilinear} $\partial_{x}$
lowers to (the first-derivative analogue of \code{Laplacian}); the heat-kernel
propagator substitutes \code{GradX}$\to \ii k$ to read off a drift coefficient
$V$.

\subsection{The passes}

The IR's semantics are a sequence of tree-walking passes, each one a single
clearly-scoped transformation.

\paragraph{Linearity (always valid).}
\code{apply\_linearity(expr, fields, coords=('x','y','z'))}
(\file{spatial\_operator\_ir.py:147}) pushes every operator node through sums and
constant coefficients: $\Lap(a\,\delta\phi+b\,\delta\psi)\to a\,\Lap(\delta\phi)+b\,\Lap(\delta\psi)$
for field- and coordinate-independent $a,b$. The inner helper \code{\_lin\_node}
first \code{.expand()}s the argument so binomial powers become sums (so
$\Lap((\bar\phi+\delta\phi)^{3})$ distributes term by term, the Cahn--Hilliard
$\nabla^{2}\phi^{3}$ case), distributes over $+$, pulls out the constant factor,
and rebuilds the operator on the fluctuation part. An argument with no
fluctuation left --- $\Lap(\bar\phi)$, $\Lap(\bar\phi^{2})$, or a pure coordinate
coefficient --- stays \emph{atomic}.

\begin{gotcha}
A factor that depends on a \emph{coordinate} is deliberately \emph{not} pulled
out: $f(x)\,\nabla^{2}\phi$ is left atomic, because its Leibniz/product rule and
the resulting $\hat f(p)$ momentum injection are a separate, deferred layer
(\file{spatial\_operator\_ir.py:51-53}). \code{apply\_linearity} is the
\emph{constant}-coefficient rule only.
\end{gotcha}

\paragraph{Saddle expansion (mean kept).}
\code{expand\_about\_saddle(expr, replacements, ...)}
(\file{spatial\_operator\_ir.py:220}) substitutes
$\phi\to\bar\phi+\delta\phi$ per \code{replacements = \{field: (mean, fluct)\}}
and re-applies linearity, so $\Lap(\bar\phi+\delta\phi)\to\Lap(\bar\phi)+\Lap(\delta\phi)$.
The mean term is \emph{retained}.

\paragraph{Contingent annihilation.}
\code{kill\_means(expr, mean\_syms, ops=('Lap','Dt','Dx'))}
(\file{spatial\_operator\_ir.py:245}) sends $\mathrm{Op}(\text{arg})\to0$
whenever \code{arg} depends only on the mean symbols. This is the
\emph{contingent} fact that the saddle is spatially homogeneous (for
$\Lap$/$\partial_{x}$) and/or stationary (for $\partial_{t}$).

\begin{gotcha}
\code{apply\_linearity} (always-valid algebra) and \code{kill\_means} (contingent
on a homogeneous saddle) are \emph{two separate passes} on purpose. For an
\emph{inhomogeneous} saddle --- a front, a pattern --- you omit \code{Lap} from
\code{ops} so that $\Lap(\bar\phi)$ \emph{survives}; it cancels the rest of the
stationarity condition rather than being silently dropped
(\file{spatial\_operator\_ir.py:36-41}). Folding annihilation into linearity
would quietly destroy this.
\end{gotcha}

\paragraph{Lowering to ring generators.}
\code{to\_derived\_generators(expr, fluct\_syms, prefix='Dg')}
(\file{spatial\_operator\_ir.py:275}) is the trick that lets the existing
multivariate-Taylor \code{expand} machinery work unchanged. It replaces each
atomic $\mathrm{Op}(\delta\phi)$ --- e.g. $\Lap(\delta\phi)$,
$\Lap(\Lap(\delta\phi))$, $\partial_{x}(\delta\phi)$ --- with a fresh ring
generator symbol, so Taylor expansion treats $\nabla^{2}\delta\phi$ as if it
were just another field (the ``$u=\delta\phi,\ v=\nabla^{2}\delta\phi$'' trick).
It returns \code{(expr2, genmap)} where \code{genmap[gen] = (base, op\_chain)}
and \code{op\_chain} records the operators applied. It works bottom-up: an
operator wrapping an already-introduced generator extends that generator's chain,
so $\Lap(\Lap(\delta\phi))$ resolves to a single $\nabla^{4}$ generator with
chain \code{(('Lap',),('Lap',))}.

\begin{defn}
A \term{derived generator} is a fresh polynomial-ring symbol standing for a
derivative of a field. Treating $\nabla^{2}\delta\phi$ as an independent
generator is what lets the framework's ordinary commutative-ring Taylor
expansion handle a derivative vertex with \emph{zero} special-casing --- the
derivative information is parked in the \code{genmap} and reapplied later by the
Fourier rules.
\end{defn}

\paragraph{Composing the passes.}
\code{prepare\_action(S, fields, replacements=None, homogeneous=True, ...)}
(\file{spatial\_operator\_ir.py:346}) ties them together for two authoring
conventions: if \code{replacements=None} the action is already in fluctuation
fields, so it only applies linearity and lowers; if \code{replacements} is given
it expands about the saddle, kills means (for a homogeneous saddle), then lowers.

\paragraph{Classification.}
\code{classify\_generators(expr, genmap, fluct\_syms)}
(\file{spatial\_operator\_ir.py:388}) splits the generators into \emph{bilinear}
(appearing only in field-degree-$\le2$ terms, so they fold into the propagator
kernel $K(\omega,k)$) and \emph{derivative-vertex} (appearing in some
field-degree-$\ge3$ term, so they carry per-leg form factors). A generator's
degree contribution is the field-degree of its base. The worked distinction in
the docstring: KPZ's $\tilde\phi\,(\partial_{x}\delta\phi)^{2}$ is degree
$1+1+1=3$, a vertex; reaction-diffusion's $\tilde\phi\,\nabla^{2}\delta\phi$ is
degree $2$, bilinear.

\subsection{The Fourier rules and form factors}

\code{form\_factor(chain, k, omega=None)} (\file{spatial\_operator\_ir.py:420})
is the Fourier image of an operator chain on a leg of momentum $k$, frequency
$\omega$:
\begin{equation}
  \Lap\ \to\ -|k|^{2},\qquad
  \partial_{t}\ \to\ -\ii\omega,\qquad
  \partial_{x_{i}}\ \to\ \ii k_{i},
  \label{eq:fourierrules}
\end{equation}
composed multiplicatively along the chain (so $\nabla^{4}=\Lap\circ\Lap\to|k|^{4}$).
The code is a direct transcription:

\begin{lstlisting}
    for entry in chain:
        op = entry[0]
        if op == 'Lap':
            out *= -k2
        elif op == 'Dt':
            ...
            out *= -I * SR(omega)
        elif op == 'Dx':
            out *= I * kv[int(entry[1])]
        elif op == '__mean__':
            out *= SR(0)                 # an un-annihilated mean derivative
\end{lstlisting}

A derivative \emph{bilinear} term (degree $\le2$) folds into the propagator
kernel via these rules. A derivative \emph{vertex} (degree $\ge3$) instead
deposits a per-leg \term{form factor} $\mathcal{R}(q,\ell)$ on the diagram, a
product over the interaction vertices:
\begin{equation}
  \mathcal{R}(q,\ell)=\prod_{\text{vertices }v}\mathfrak{F}(v),
  \qquad
  \mathfrak{F}(v)=\sum_{t:\,n_{\text{phys}}(t)=\deg(v)}w_{t}\,\mathfrak{f}_{t}(v),
  \label{eq:Rcal}
\end{equation}
where $w_{t}=c_{t}/\sum c$ is the coupling weight and $\mathfrak{f}_{t}$ is the
type's Fourier kernel, evaluated either at the response-leg momentum
(\term{composite} mode: $\nabla^{2}(\phi^{2})$, $\partial_{x}(\phi^{2})$) or as a
product over physical-leg momenta (\term{perleg} mode: $(\partial_{x}\phi)^{2}$).

\begin{defn}
\term{Composite} vs \term{perleg} is the distinction between an operator acting
on the whole $\phi^{n}$ composite at a vertex (one factor, the response-leg
momentum) versus acting on each physical leg separately (a product of per-leg
factors). $\nabla^{2}(\phi^{2})$ is composite; $(\partial_{x}\phi)^{2}$ is
perleg. The vertex-terms table records each type's \code{mode}, and the diagram
form-factor assembler dispatches on it.
\end{defn}

The loop integrator must then average this polynomial form factor over the
loop-momentum Gaussian. There are two routes: a numerical Gauss--Hermite
quadrature, or --- the principled route --- a closed-form joint $(\ell,q)$
Gaussian moment, described next.

\section{The analytic spatial inverse Fourier transform}

A derivative vertex makes the loop-momentum integrand a Gaussian times a
polynomial form factor, and the polynomial part can be integrated \emph{exactly}
by Wick's theorem rather than sampled numerically. This retires the
$q$-grid entirely for $d=1$ derivative-vertex theories.

\code{\_build\_wick\_moment(Rcal, ls, qs)} (\file{pipeline\_bridge.py:815}) is
the construction (it is the $k=2$, $d=1$ case; \code{\_build\_wick\_moment\_multi}
at line~1893 generalizes to $k\ge3$). It splits the loop momentum
$\ell=a\cdot q+\xi$, turns the external $q$ into a complex Gaussian via the
Fourier-transform source, and writes the per-chamber inverse FT as
$\delta C(x)=A\cdot K(\mathcal{B},x)\cdot M_{F}$ with
\begin{equation}
  M_{F}=\mathbb{E}_{Q\sim\mathcal{N}(c,s)}\,
        \mathbb{E}_{\xi\sim\mathcal{N}(0,\Sigma)}
        \big[\mathcal{R}(a\,Q+\xi,\,Q)\big],
  \label{eq:MF}
\end{equation}
both expectations being closed-form Gaussian moments --- a non-central moment
$\mathbb{E}[Q^{m}]$ for the external part and an \term{Isserlis} pairing sum for
the loop part.

\begin{defn}
The \term{Isserlis (Wick) theorem} states that the expectation of a product of
jointly-Gaussian zero-mean variables equals the sum, over all perfect matchings
of the variables, of the products of pairwise covariances; an odd number of
factors gives zero. \code{\_isserlis(idx, Sget)} (\file{pipeline\_bridge.py:802})
is the recursive symbolic implementation: a sum over all perfect matchings of the
index multiset of $\prod_{\text{pairs}}\Sigma$. This is what makes the
loop-momentum average a finite \emph{exact} sum rather than a quadrature.
\end{defn}

The construction is built once per diagram and then \term{lambdified} into NumPy
callables. The key performance factorization, recorded in the docstring of
\code{moment\_x}: the polynomial-in-$X$ form factor has coefficients $g_{k}$ that
are \emph{$x$-independent}, so the expensive per-sample $g_{k}$ are evaluated once
and contracted with cheap powers of $X$ --- roughly an $n_{x}$ speedup over
recomputing the moment at every output point. A second variant,
\code{moment\_bessel}, grades the form factor by the radial scaling power for the
Bessel-K backend.

\begin{defn}
\code{sympy.lambdify(args, expr, 'numpy', cse=True)} \emph{compiles} a SymPy
symbolic expression into a fast Python function whose body is NumPy operations,
so a whole array of samples can be evaluated at once. \code{cse=True} runs
\term{common-subexpression elimination} first --- factoring out shared subterms
like $1/\mathcal{B}^{k}$ --- which is a real speedup for the Wick-moment
polynomials. This is the single most-used bridge from symbolic algebra to
numerics in the whole subsystem (\file{pipeline\_bridge.py:911,937,997,1043,1979}).
\end{defn}

\begin{note}
SymPy and Sage are \emph{both} present in this subsystem, deliberately. Sage's
Symbolic Ring \code{SR} is the canonical physics-facing algebra --- the action,
the inverse propagator $K_{\mathrm{ft}}$, the parameter substitutions. SymPy
(imported lazily as \code{import sympy as \_sp} inside functions) handles the
momentum-space \emph{polynomial} bookkeeping of diagrams --- routing momenta,
building form factors, taking Wick moments --- because that work is polynomial
algebra that does not need Sage's heavier algebraic-number machinery. They never
mix in a single expression.
\end{note}

\section{The bridge architecture and the dispatch tree}\label{sec:dispatch}

We can now assemble the pieces into the dispatcher,
\file{pipeline\_bridge.py}. Its defining trick is the \term{Laplacian
substitution}.

\subsection{The Laplacian substitution and certification}

\code{pipeline\_C\_q\_tau} (\file{pipeline\_bridge.py:259}) runs the shared,
time-only diagram pipeline at $\Lap=-q^{2}$:

\begin{lstlisting}
    Lap = SR.var('Laplacian')
    nps = dict(base_np_sr)
    nps[Lap] = -(qval ** 2)
    compute_poles_and_residues(prop, nps, verbose=False)
    ...
    res = compute_correction_td(typed_diagrams=[r[0] for r in records],
                                prefactors=[r[1] for r in records],
                                k=k, propagator_data=pdata, ...)
\end{lstlisting}

The line \code{nps[Lap] = -(qval**2)} injects $\Lap\to-q^{2}$ so the time-only
pipeline sees a rational propagator at effective mass $m(q)=\mu+Dq^{2}$, and
\code{compute\_poles\_and\_residues} re-solves the pole/residue cache per $q$,
exactly as \code{compute.py} does for a temporal model. The diagram answer
$C(q,\tau)$ is then handed to \code{certify\_modes}
(\file{pipeline\_bridge.py:286}), which computes the maximum relative error
between the pipeline's $C(q,\tau)$ and the per-mode reference
$\sum_{\alpha}\kappa_{\alpha}/(\mu_{\alpha}+D_{\alpha}q^{2})\,
\ee^{-(\mu_{\alpha}+D_{\alpha}q^{2})|\tau|}$ over a small sample grid.

\begin{note}
This is the operational meaning of ``route through the pipeline, then
certify.'' The Laplacian substitution lets the engine reuse the entire temporal
diagram machinery --- enumeration, classification, time-domain integration ---
with momentum carried as a parameter. Certification then confirms that the
diagrams the pipeline produced \emph{are} the heat-kernel modes the analytic
$q\to x$ transform is about to use. For a single-field reaction-diffusion theory
the relative error is around $10^{-15}$ at tree level.
\end{note}

\begin{gotcha}
Certification is a \emph{hard gate}, not a warning.
\code{compute\_spatial\_correlator\_via\_pipeline} \emph{raises}
\code{SpatialPropagatorError} if \code{certify\_max\_rel} exceeds
\code{certify\_tol} (default $10^{-8}$). The message says the $(\mu,D,\kappa)$
extraction or the diagram routing is wrong for the theory. A failed
certification is the engine refusing to return a number it cannot vouch for.
\end{gotcha}

\subsection{The five drivers and the branch points}

\code{compute\_cumulants} chooses among the bridge's drivers based on a short
sequence of predicates. The decisive ones are:

\begin{itemize}
  \item \code{prop['G\_tx\_sym'] is None} $\Rightarrow$ off-diagonal coupling
        $\Rightarrow$ the \emph{coupled} path. (The diagonal heat-kernel block is
        absent; the propagator's \code{K\_ft} is still present and the coupled
        path reads it.)
  \item \code{\_field\_index(leg0) != \_field\_index(leg1)} $\Rightarrow$ a
        \emph{cross}-correlator $\Rightarrow$ the coupled path (the diagonal path
        would wrongly reconstruct $\avg{\phi_{i}\phi_{i}}$).
  \item \code{ref.is\_scalar\_diffusion} (i.e. $\widehat{\mathcal{D}}=0$)
        $\Rightarrow$ the exact \code{coupled\_two\_point}; otherwise consume the
        Dyson policy and use \code{dressed\_tree\_C}.
  \item \code{max\_ell >= 1} $\Rightarrow$ a loop driver; if the tree was coupled,
        the coupled loop driver.
  \item \code{len(external\_fields) >= 3} $\Rightarrow$ a $k$-point driver.
\end{itemize}

\begin{figure}[ht]
\centering
\begin{tikzpicture}[
    >={Stealth[round]}, node distance=6mm,
    box/.style={draw=noteframe, fill=notebg, sharp corners, rounded corners=1pt,
                font=\footnotesize, align=center, inner sep=4pt, text width=33mm},
    leaf/.style={draw=defframe, fill=defbg, sharp corners, font=\footnotesize\ttfamily,
                 align=center, inner sep=4pt, text width=46mm},
    q/.style={font=\scriptsize\itshape, fill=white, inner sep=1pt}]
  \node[box] (start) {\code{compute\_cumulants}\\ (spatial)};
  \node[box, below=10mm of start] (coupled) {diagonal heat-kernel\\ block present?};
  \node[leaf, right=24mm of coupled] (treepipe)
        {compute\_spatial\_\\correlator\_via\_pipeline};
  \node[box, below=12mm of coupled] (scalar) {scalar diffusion\\ ($\widehat{\mathcal D}=0$)?};
  \node[leaf, right=24mm of scalar] (dyson) {dressed\_tree\_C\\ (Dyson series)};
  \node[leaf, below=10mm of scalar] (twopt) {coupled\_two\_point\\ (Lyapunov / FDT)};
  % edges
  \draw[->] (start) -- (coupled);
  \draw[->] (coupled) -- node[q,above]{yes (+ same field)} (treepipe);
  \draw[->] (coupled) -- node[q,right]{no / cross} (scalar);
  \draw[->] (scalar) -- node[q,above]{no} (dyson);
  \draw[->] (scalar) -- node[q,right]{yes} (twopt);
\end{tikzpicture}
\caption{The tree-level dispatch. A diagonal heat-kernel block and a same-field
correlator take the certified single-field path; off-diagonal coupling or a
cross-correlator routes to the coupled path, which then branches on scalar
vs.\ unequal diffusion. (\code{max\_ell}$\ge1$ and $k\ge3$ select loop and
$k$-point drivers on top of this.)}
\label{fig:dispatch}
\end{figure}

The coupled-tree driver \code{compute\_coupled\_tree\_correlator}
(\file{pipeline\_bridge.py:449}) reads \code{prop['K\_ft']}, extracts
$(M,\mathcal{D},V)$ via \code{reaction\_diffusion\_matrices}, numericizes them,
and \emph{requires $V=0$} (a genuine drift is not supported in the
scalar-diffusion v1 --- it raises \code{NotImplementedError} otherwise). It then
builds the \code{SpectralReference}, and:

\begin{itemize}
  \item if scalar diffusion, defines $C_{ij}(q,\tau)$ as
        \code{coupled\_two\_point(ref, N, q**2, tau)[i,j]};
  \item if unequal diffusion, consumes the Dyson policy
        (\code{model['spatial']['dyson']} or the \code{SPATIAL\_DYSON\_ORDER} env
        override), \emph{requires} a fixed truncation order, and checks the
        convergence ratio $\rho$ of \eqref{eq:rho}.
\end{itemize}

\begin{gotcha}
\textbf{The Dyson divergence guard.} \code{compute\_coupled\_tree\_correlator}
\emph{raises} \code{SpatialPropagatorError} if $\rho=\lVert\widehat{\mathcal{D}}\rVert/D_{0}\ge1$
(``the large-$|k|$ insertion outgrows the reference'') and warns if $\rho>0.6$
(slow convergence). The series is exact at $q=0$ regardless --- the insertion
vanishes there --- but for the whole grid to converge you need $\rho<1$. Retune
$D_{0}$ to shrink it (\file{pipeline\_bridge.py:519-523}).
\end{gotcha}

For $d=1$ infinite boundary with scalar diffusion the external $q\to x$ transform
is done by the \term{analytic spectral inverse FT}: each $(\alpha,\beta)$ term
$(P_{\alpha}NP_{\beta}^{\mathsf T})/(m_{\alpha}+m_{\beta}+2D_{0}q^{2})$ is a
single-mode correlator at denominator mass $\mu_{d}=(m_{\alpha}+m_{\beta})/2$,
transformed exactly via \code{free\_two\_point}$(\mu_{d},D_{0},\tfrac12;x,\tau)$,
with $C_{ij}(-\tau)=C_{ji}(\tau)$.

\section{The single-field generic loop path}

The single-field loop driver \code{compute\_spatial\_correlator\_generic}
(\file{pipeline\_bridge.py:1434}) is the one genuine momentum-integral path, and
it is where the form factors of the operator IR and the analytic IFT of the
previous sections actually get used. Its outline:

\begin{enumerate}
  \item Compute the tree $C_{0}$ via the certified
        \code{compute\_spatial\_correlator\_via\_pipeline}; if the tree turned
        out coupled, route to \code{compute\_coupled\_loop\_correlator}.
  \item Build the numeric derivative-vertex form-factor table from
        \code{ns.\_operator\_ir\_vertex\_terms} (substitute couplings into the
        symbolic weights $c_{t}/\sum c$).
  \item Enumerate records, map every diagram to a (descriptor, prefactor,
        form-factor callable, $\ell$) tuple, and filter to the live ones.
  \item Choose an integrator backend, check the memory guard, and do either the
        analytic or the numerical $q\to x$ transform.
\end{enumerate}

\begin{gotcha}
\textbf{Drift is refused on this path too.} For any field with a non-zero saddle
drift $V\ne0$ (a genuine advection $v\cdot\partial_{x}\phi$), the driver raises:
the drifting propagator is validated only at the heat-kernel oracle level, not
in the momentum integrator (which uses $m_{k}=\mu+Dk^{2}$). Gradient theories
like Burgers and KPZ, where $V\to0$ at $\phi^{*}=0$, run fine
(\file{pipeline\_bridge.py:1540-1554}).
\end{gotcha}

\subsection{Backends, the memory guard, and the analytic/numerical gate}

The \code{SPATIAL\_INTEGRATOR} environment variable selects the backend: the
default \code{grid} (deterministic product quadrature), \code{mc}
(importance-sampled Monte-Carlo, bounded memory but biased for derivative
vertices), or \code{bessel} (radial Bessel-K times angular Monte-Carlo, which
regularizes the $\det M\to0$ singularity and handles derivative vertices at two
loops).

\begin{gotcha}
\textbf{The memory guard / curse of dimensionality.} A chamber's quadrature is
$P=n_{t}^{\,n_{V}}\cdot n_{s}^{\,n_{C}}$ samples ($n_{V}$ internal vertices,
$n_{C}$ correlation edges). A KPZ two-loop diagram has $n_{V}=4$, $n_{C}=3$,
giving $P\approx1.8\times10^{8}$ at $(n_{t}{=}16,n_{s}{=}14)$, and the
$(P,n_{x})$ heat-kernel array alone is tens of gigabytes --- an out-of-memory
crash. The guard (\file{pipeline\_bridge.py:1633-1665}) refuses \emph{up front}
with the exact numbers and four mitigation knobs: lower \code{max\_ell}, use the
\code{mc} backend, coarsen \code{SPATIAL\_GRID\_NT}/\code{SPATIAL\_GRID\_NS}, or
raise \code{SPATIAL\_MEM\_BUDGET\_GB} (default 6). It is bypassed for the
bounded-memory \code{mc}/\code{bessel} backends.
\end{gotcha}

The final branch is the \term{analytic-vs-numerical FT gate}
(\file{pipeline\_bridge.py:1677}):
\begin{lstlisting}
    _all_plain = all(rec[2] is None for rec in live_g)
    _force_num = _os.environ.get('SPATIAL_FORCE_NUMERICAL_FT', '') == '1'
    _use_analytic = (_all_plain or d == 1) and not _force_num
\end{lstlisting}
The analytic path (Chapter~\ref{ch:spatial-heatkernel} for plain vertices; the
Wick-moment construction above for $d=1$ derivative vertices) computes
$\delta C(x,\tau)$ directly --- no $q$-grid, no ringing. The numerical path
computes $\delta C(q,\tau)$ on a grid and then Fourier-transforms; it is reachable
as a cross-check via \code{SPATIAL\_FORCE\_NUMERICAL\_FT=1} but is exact only in
the $q_{\text{cut}}\to\infty$, $n_{q}\to\infty$ limit.

\begin{gotcha}
\textbf{The \texttt{Rcal == 0} form factor must still evaluate to zero on every
path.} A diagram that vanishes by a conservation law (e.g. the
$\nabla^{2}(\phi^{2})$ tadpole with a forced-zero leg momentum) returns
\code{\_zero\_ff\_with\_moments}, which attaches zero-returning \code{moment\_x},
\code{moment\_x\_multi}, and \code{moment\_bessel} callables --- otherwise the
analytic IFT would hit a missing-moment error instead of contributing zero
(\file{pipeline\_bridge.py:972-990}).
\end{gotcha}

\subsection{The bubble-versus-tadpole test}

One small but load-bearing classifier deserves mention.
\code{\_diagram\_is\_bubble(td)} (\file{pipeline\_bridge.py:664}) decides whether
a one-loop diagram is a momentum-\emph{dependent} bubble:

\begin{lstlisting}
    for v in route_momenta(td).edge_k2().values():
        e = _sp.expand(v)
        syms = sorted(e.free_symbols, key=str)
        for ii in range(len(syms)):
            for jj in range(ii + 1, len(syms)):
                if _sp.expand(e.diff(syms[ii]).diff(syms[jj])) != 0:
                    return True
    return False
\end{lstlisting}

It routes the momenta, then for each edge's $k^{2}$ checks for a nonzero
\emph{mixed second partial} --- a $q\cdot\ell$ cross term. A bubble has such a
term (some edge carries $(q-\ell)^{2}$); a tadpole has every edge at pure
$q^{2}$, $\ell^{2}$, or $0$. The test is topology-agnostic.

\begin{gotcha}
A topological bubble whose scalar prefactor is $\propto\phi^{*2}$ is \emph{dead}
at $\phi^{*}=0$ and must not trigger the bubble route.
\code{\_prefactor\_is\_live} (\file{pipeline\_bridge.py:1058}) substitutes the
saddle-carrying \code{num\_params} and filters these out: only \emph{live}
bubbles count. (If free symbols remain it conservatively returns true.)
\end{gotcha}

\section{General-\texorpdfstring{$k$}{k} cumulants}\label{sec:kpoint}

For correlators with three or more external legs the bridge has two more drivers,
the $k$-point companions of the tree and loop paths.
\code{compute\_spatial\_kpoint} (\file{pipeline\_bridge.py:1788}) is the
single-field general-$k$ path; \code{compute\_coupled\_kpoint}
(\file{pipeline\_bridge.py:2000}) is the coupled scalar-diffusion companion. Both
take \code{points}, an array of $(x_{j},\tau_{j})$ offsets of slots
$1,\dots,k-1$ relative to slot~0.

The key structural difference from $k=2$ is how external legs are wired. The
$k$-point drivers use an \term{orbit-stabilizer external-Wick mapping sum}
(\code{field\_respecting\_mappings} plus \code{external\_wick\_compensation})
rather than the explicit orientation and $\pm\tau$ completion of the two-point
case. As the comment at \file{pipeline\_bridge.py:509} notes, the two views
coincide at $k=2$; at higher $k$ the mapping sum is the general machinery. Points
that share a $\tau$-configuration are batched, and the coupled $k$-point driver
is scoped to scalar diffusion, plain vertices, and infinite boundary.

\section{Operational guardrails}

Several of the boxes above are instances of a single design stance: this
subsystem prefers to \emph{refuse} a computation it cannot certify rather than
return a plausible-looking wrong number. Collected in one place, the guardrails
are:

\begin{itemize}
  \item \textbf{The macOS fork prohibition (the big one).} Fork-based parallelism
        was \emph{removed} from the spatial path. Forking a Jupyter kernel on
        macOS after the GUI/BLAS libraries have initialized crashes the kernel
        \emph{and} the OS, even with a single worker. The only safe speedup is
        thread-based (NumPy releases the GIL during heavy array ops) or
        $q$-grid NumPy batching. The numerical-FT path uses a
        \emph{smart-gated} \code{ThreadPoolExecutor}
        (\file{pipeline\_bridge.py:1713}): threads engage only when a
        loop-order-$\ge2$ diagram is present (about $2.5\times$ at two loops;
        \emph{slower} at one loop, where the tiny matrices are dispatch-bound, so
        one loop stays serial), with accumulation on the main thread for
        bit-identical results.
  \item \textbf{The memory budget.} As above: refuse up front when a chamber's
        quadrature would exceed \code{SPATIAL\_MEM\_BUDGET\_GB}.
  \item \textbf{Certification is a hard gate.} A failed mode certification raises
        \code{SpatialPropagatorError}.
  \item \textbf{Defective reaction matrix.} \code{spectral\_projectors} raises if
        the eigenvector condition number exceeds $10^{10}$.
  \item \textbf{Dyson divergence.} $\rho\ge1$ raises; $\rho>0.6$ warns.
  \item \textbf{Drift $V\ne0$ is refused} on both the single-field generic and the
        coupled-tree paths.
  \item \textbf{Never default a leg index silently.} Both
        \code{\_legs\_to\_phys\_idx} (\file{pipeline\_bridge.py:89}) and the
        per-leg index resolver \emph{raise} rather than fall back --- a wrong leg
        would quietly compute a \emph{different} correlator ($C_{aa}$ instead of
        $C_{ab}$). The point is emphasized twice in the source comments
        (\file{pipeline\_bridge.py:122-124,542-543}).
\end{itemize}

\begin{defn}
The \term{GIL} (Global Interpreter Lock) is the CPython mechanism that lets only
one thread execute Python bytecode at a time. The reason \emph{thread}
parallelism can still help here is that NumPy's heavy array operations release
the GIL while they run in compiled C --- so several threads can do NumPy work
genuinely in parallel. Process/fork parallelism would side-step the GIL too, but
it is forbidden here for the macOS-crash reason above.
\end{defn}

\begin{note}
\textbf{Environment escape hatches.} Several knobs exist for cross-checks and
debugging, not normal use: \code{SPATIAL\_FORCE\_NUMERICAL\_FT=1} keeps the
numerical-FT path reachable; \code{SPATIAL\_Q\_CUT} and \code{SPATIAL\_N\_Q} tune
it; \code{SPATIAL\_DYSON\_ORDER} overrides the model's Dyson policy;
\code{SPATIAL\_INTEGRATOR}, \code{SPATIAL\_GRID\_NT}/\code{NS},
\code{SPATIAL\_MC\_N}, and \code{SPATIAL\_MEM\_BUDGET\_GB} tune the loop
integrator. In ordinary use you set none of them.
\end{note}

\section{The oracle module: \texttt{loop\_dyson.py}}

The fifth module, \file{loop\_dyson.py}, carries a prominent banner
(\file{loop\_dyson.py:1-5}):

\begin{quote}\itshape
``$\triangle$ ORACLE-ONLY --- not on the production path. Superseded by
\code{full\_integrator.py}; reached only by its own test(s).
\code{compute\_cumulants} does NOT use this module.''
\end{quote}

It is an \emph{independent} implementation of the one-loop $\tilde\phi\phi^{2}$
reaction-diffusion \term{bubble} self-energy via the MSR Dyson equation --- both
the equal-time structure-factor correction $\delta C(q,0)$ and the full
$\tau$-dependent correction. It computes the retarded and Keldysh bubbles
($\Sigma_{R}=\int\,\dd\ell\,F(\ell)\,\Gret(\ell,t)\,C(q-\ell,t)$ and the
even-in-$t$ Keldysh analogue) by adaptive Gauss--Kronrod quadrature
(\code{scipy.integrate.quad}), and assembles the Dyson terms with the pinned
normalization \code{C\_R, C\_K = 4.0, 2.0} (\file{loop\_dyson.py:89}), validated
to give a fitted coefficient $B=0.99$ against a perturbative simulation.

Its value is precisely that it shares \emph{no code} with the production
\code{full\_integrator} path: a number that both produce independently is a
genuine cross-check. But the $C_{R}=4,C_{K}=2$ normalization and the $B=0.99$
validation are historical context, not the live computation.

\begin{defn}
\code{scipy.integrate.quad(f, a, b, limit=...)} is \term{adaptive Gauss--Kronrod
quadrature} (QUADPACK's \code{QAGS}/\code{QAGI}). Unlike the fixed-node
Gauss--Legendre rule, it recursively subdivides the interval until a local error
estimate is met; \code{limit} caps the number of subdivisions. The oracle uses it
for the self-energy momentum integrals and the Dyson time integrals.
\end{defn}

\section{What you just learned}

This chapter covered the part of the engine that handles every spatial case the
diagonal path cannot:

\begin{enumerate}
  \item \textbf{Coupled fields} need a matrix propagator. The reference-diffusion
        split $\mathcal{D}=D_{0}I+\widehat{\mathcal{D}}$ and the spectral
        resolution $M=\sum m_{\alpha}P_{\alpha}$ give a closed-form reference
        $G_{0}$ that is \emph{exact} when diffusion is scalar
        (\file{spectral\_propagator.py}).
  \item \textbf{The tree two-point} of a coupled linear theory is the
        Lyapunov/fluctuation--regression $C(q,\tau)=\ee^{-A|\tau|}\Sigma$, not a
        sum of independent modes (\code{coupled\_two\_point}).
  \item \textbf{Unequal diffusion} is handled by the Dyson--Duhamel series in
        powers of $\widehat{\mathcal{D}}|k|^{2}$, whose one new primitive
        $\Phi_{n}$ is evaluated confluent-safely by the Opitz matrix exponential
        (\file{dyson\_dressing.py}).
  \item \textbf{Derivative vertices} are hosted as inert operator nodes in an IR
        with explicit linearity, saddle-expansion, annihilation, and lowering
        passes; each deposits a momentum form factor whose loop average is taken
        either numerically or exactly by the Wick-moment construction
        (\file{spatial\_operator\_ir.py}, \code{\_build\_wick\_moment}).
  \item \textbf{The bridge} ties it all to the shared diagram pipeline via the
        $\Lap\to-q^{2}$ substitution, certifies the diagram answer against the
        per-mode reference, dispatches to the right driver, and guards
        aggressively against the cases it cannot certify
        (\file{pipeline\_bridge.py}).
\end{enumerate}

The recurring lesson is the one in the bridge's docstring and in nearly every
\emph{Gotcha} box above: the engine reuses the temporal machinery wherever it
can, and where it cannot prove the answer is right it refuses to return one. The
next part of the manual steps back out to the assembly layer ---
\code{compute\_cumulants} itself and the user-facing engine API --- which is what
selects every driver this chapter described.
